
\section{The Congregation and Lay Activity}

Editorial article in “Folkebladet” for April 10, 1907.

In the previous number of “Folkebladet” we had to write something about the manner in which the great association committee treated “lay activity.” It is rather unpleasant to have to take up such misleading representations of a precious cause; and therefore this time we owe it to ourselves to speak more positively about the matter.

When, then, we wish to speak about the congregation and lay activity, it is in order to say that it is to be expected that there will be some difference for the cause of lay activity as well, where there is a congregation.

By congregation we mean a gathering of active Christians, who use the means of grace and the gifts of grace for their own and their fellow men’s salvation and blessedness. Where such a gathering of active Christians is found, the question does not become one of binding and controlling lay activity and other gifts of grace, but much more of giving them room and freedom, occasion and encouragement, so that they may come to the greatest flourishing and the most effective use.

Here there arises, then, a question of which there was no occasion to speak in the previous article, namely: who has this gift of speaking in an assembly, which is indeed one of the conditions for lay activity? And who shall decide whether such a gift is present? And how is such a gift best discovered and developed?

Here lies a difficulty which all have felt who with interest and love have thought about lay activity and its flourishing. Not everyone has received the gift because he imagines that he has it. Nor does anyone receive the gift by being called by a congregation to speak in an assembly. It is seen, indeed, with some of the pastors, that they themselves believe they have the gift, and they also have the congregation’s call, but nevertheless the gift of grace does not seem to be present.

With regard to this matter, we have heard that Hauge’s friends had the custom that some of the old, tried Christians and proclaimers of God’s Word exercised the task of testing spirits and gifts, and that new and younger forces had to submit to passing through a fairly long time of testing before they were given an independent place in the ranks of the witnesses. And a similar procedure should no doubt be sought also among us. We say “similar,” for it is probably most difficult to attain altogether the same manner.

It is clear that neither chairs nor examination can produce a lay preacher; but neither does anyone become a lay preacher by simply ceasing from other work and making it his livelihood to travel about and preach. A traveling evangelist or revivalist, who has this work as his only occupation and his livelihood, is not a lay preacher in the older sense of the word. We understand also that those who make the preaching task their only work are preferably called emissaries both in Norway and here among us.

Without doubt the work of an evangelist has its great significance. But it is not the Haugean lay activity. We do not depreciate the one work because we write in commendation of the other.

A lay preacher is a preacher who has a message from God to bring to his fellow men, and who cannot, without unfaithfulness, be silent about the message which God has given him to bear forth; and who therefore both in conversation with individuals and in greater or smaller assemblies of various kinds bears forth this message. But he remains a layman while thus preaching, and does not have his livelihood from this work of preaching.

They are driven to speak and cannot refrain from it; they begin where they have opportunity, in their home and their home district; they travel about for shorter or longer distances, and neither congregational resolutions nor other prohibitions bar their way.

Of course they are dependent upon the trust they gain, especially among other Christians. To give them a congregational call is not necessary and therefore not desirable either; for then they could not order either their time or their speech according to their ability and gift, but would have to try to please those in whose pay they were. But with that it is not excluded that the congregation may encourage and urge them forward in such manner as it itself finds right and useful.

But their gift is tested by the use of the gift and proved by the hearing they gain. There are honorable and upright ways by which their influence may gradually be extended to wider circles, if they attain the glorious thing of gaining the Christians’ trust. But to procure for oneself a whole number of attestations, which perhaps are given half unwillingly because no one wishes to offend them by refusing them an attestation, is often tantamount to a lack of real trust.

The true lay preacher is independent of gifts of money. He has his temporal work and his livelihood apart from his Christian testimony. Therefore his authority is greater and his testimony often more effective among many people.

Such men stand as free and independent as is at all compatible with the Christian state and brotherhood. And therein lies their value for the congregation in no small degree. They belong to the congregation, and they belong to the Lord; they are driven by the Lord’s Spirit to bear witness for Him in the congregation and for those outside the congregation. The congregation cannot, without overreach, command them to be silent, when it must be acknowledged that God has given them message and gift; nor can they demand of the congregation that it should support them, when it has not called them. But it may be expected, where there is a free and living congregation, that there will be much such lay activity, and that broad room and free opportunity will be opened to it. For though not many in each place have great gifts in this direction, yet many Christians ought to have Spirit and strength to bear such witness as may be for the edification of those who hear it.

We are altogether convinced that living congregations will do well to heed the apostolic rule in 1 Cor. 14:

“What then is to be done, brothers? When you come together, each one has a psalm, he has a teaching, he has an utterance, he has an interpretation; let all things be done unto edification!”

And it is surely not too much to say that if, according to this apostolic rule, all Christians become the Savior’s witnesses also “in an assembly,” lay activity too will flourish better. For in the first place the gifts come into use in the right Christian way, and in the second place a lay preacher does not come to stand so alone as a “lone tree” on the prairie, which is always a dangerous and distressful position.

Let a little life come into the work, and it will go better with us in every direction!

Georg Sverdrup
